\section{La escala no es un número, sino varias colas que no pueden mezclarse}

La sección anterior estableció la cadena de respuesta; esta sección explica por qué el «crecimiento de la escala» desborda a las organizaciones. Las estadísticas públicas suelen presentar al mismo tiempo reports, files, accounts, incidents y victims, cuyas unidades, tasas de duplicación y criterios anuales difieren. Si un equipo de ingeniería escribe en el dashboard solo la cifra más grande, estimará mal la carga de almacenamiento, revisión, investigación y entrenamiento de modelos.

\subsection{Reports, files, duplicates y novel candidates}

Un report puede contener varios archivos, y un mismo archivo puede ser reportado varias veces por distintas plataformas o cuentas; lo que los investigators realmente necesitan atender con prioridad suele ser la parte que puede corresponder a una nueva víctima, a un riesgo persistente o a material antes desconocido. La \term{revictimization} (nuevo daño) se refiere al daño repetido que sufre la víctima cuando el material se copia, se ve o se difunde de forma continua; por ello, incluso la reaparición de un «archivo conocido» merece bloquearse con rapidez, pero no debe confundirse con el «número de casos nuevos».

\lecturefigure{02-scale-funnel.png}{El volumen de reportes, el volumen de archivos, el contenido duplicado, los candidatos desconocidos y la escasa capacidad de revisión son medidas distintas.}{Subtítulos locales 00:00--05:50; nota sobre los criterios anuales de NCMEC; redibujo conceptual.}

\begin{knowledgebox}{Lectura de la figura: primero pregunte por el denominador y las reglas de deduplicación}
Al leer cualquier dato de escala, primero escriba la unidad: un report es un envío, un file es un adjunto, un candidate es un objeto preseleccionado por un algoritmo y pendiente de revisión, y un case es la unidad de investigación que una institución forma al agregar elementos según sus reglas. También hay que preguntar si se deduplica por cryptographic hash, perceptual hash, account, incidente o victim. NCMEC ha ajustado sus métodos de estadística y de reporte en distintos años, por lo que las páginas de 2023, 2024 y 2025 solo pueden interpretarse dentro de sus propias definiciones, y un punto de inflexión en la curva no puede atribuirse directamente a un cambio interanual del daño real.
\end{knowledgebox}

\begin{importantbox}{La planeación de capacidad debe considerar cuatro tasas a la vez}
Sea $\lambda_r$ el número de reportes que llegan por unidad de tiempo, $m_f$ el número promedio de archivos por reporte, $p_k$ la proporción de contenido conocido y $p_u$ la proporción de candidatos desconocidos que requieren revisión humana; entonces la tasa aproximada de llegada a revisión es
\[
\lambda_h \approx \lambda_r m_f p_u.
\]
Aquí $\lambda_h$ no es el número de víctimas ni el número de casos legales; es solo una estimación de la carga de la cola de revisión. Las coincidencias con contenido conocido siguen activando acciones de atención y reportes, por lo que el hecho de que desaparezcan de la cola humana no significa que no tengan costo.
\end{importantbox}

\teachervoice{En clase se distinguió repetidamente entre la «detección y reporte» de las plataformas y la «identificación de víctimas y persecución de los responsables» de los investigadores. El ponente enfatizó que lo que menos necesitan las instituciones al final de la cadena es otro montón de archivos sin priorizar, sin contexto y con gran cantidad de duplicados.}

\subsection{Resumen de la sección}

En esencia, el problema de escala es de colas heterogéneas: la difusión duplicada debe bloquearse a gran velocidad, el contenido desconocido debe revisarse con cuidado, las pistas de alta prioridad deben llevar contexto y los investigadores necesitan un case package accionable. Una sola cifra total no puede orientar el diseño del sistema.
